\chapter[案例研究 4：社会模拟]{案例研究 4：\\利用大语言模型智能体模拟社会}\label{chap:generative_agents}


\begin{learningbox}

在本案例研究中，你将通过广泛采用的生成式智能体（Generative Agents）框架，熟悉基于智能体的社会模拟（agent-based social simulation）\citep{park2023generative}。
具体而言，你将能够：

\begin{itemize}[leftmargin=*, nosep]
    \item 运行并配置由大语言模型（large language model，LLM）驱动的{\em 多智能体社会模拟}。
    \item 开展结构化的{\em 机器心理学（machine-psychological）}实验。
    \item 比较两种{\em 认知架构（cognitive architecture）}：标准记忆系统，以及具有负性偏向（negativity bias）的记忆系统。
    \item 探究{\em 性格描述（personality description）}与架构选择如何塑造智能体的涌现行为。
    \item 批判性地思考{\em 提示脆弱性（prompt fragility）}、模型依赖性，以及大语言模型驱动的模拟所存在的局限。
\end{itemize}

\end{learningbox}


\section{社会模拟}
在社会模拟中，生成式智能体是由大语言模型驱动的角色，它们既与环境交互，也彼此交互。智能体可以规划活动、维持记忆，或反思自身经历。\citet{park2023generative} 提出的生成式智能体框架表明，大语言模型驱动的角色聚在一起，能够产生可信的日常作息、社会交互，甚至协调一致的群体活动。这一切都涌现于一种相对简单而优雅的认知架构。

本案例研究的特别之处在于，你将为智能体使用\emph{两种}认知架构，在同一场景中分别运行，并比较涌现行为。一种架构平等对待所有记忆。另一种架构引入了\emph{负性偏向}：对负性经历进行更强的编码，使其更有可能被检索出来并成为反思对象。我们将看到，这种受认知心理学研究发现启发的不对称性 \citep{baumeister2001bad}，会对实验智能体产生可测量的影响。


本章旨在理解架构选择如何带来涌现后果；在这里，架构选择体现为智能体内部处理过程的细微变化。你还将看到，智能体的角色设定（persona）可能十分脆弱：性格描述中的一个词，或记忆权重的细微调整，都可能从根本上改变智能体行动、规划与反思的方式。

\subsection*{课程材料}
你可以在这里找到课程练习材料： \\
\url{https://github.com/DorotaMoravcikova/agentic-llm-course-excercises}




\section{背景}

\subsection{生成式智能体}

\citet{park2023generative} 提出的生成式智能体框架表明，大语言模型驱动的角色能够产生可信的日常作息、社会交互，甚至协调群体活动。

每个智能体都具有三个核心组件：
\begin{enumerate}
    \item 记忆流（memory stream）：带有时间戳的日志，记录智能体感知过、做过和想过的一切。
    \item 检索函数：当智能体需要采取行动时，根据记忆的新近性（recency）、重要性（importance），以及与当前情境的相关性（relevance），检索有关记忆。
    \item 反思机制：智能体定期从原始记忆中综合形成更高层次的认识。
\end{enumerate}

智能体的\textit{性格}由一段简短的自然语言描述定义。这段描述存放在一个 {\em scratch} 文件中，规定角色的姓名、年龄、后天习得和先天具有的特征，以及生活方式。每一条提供给大语言模型的提示都会包含这段描述，以此锚定智能体的行为。

\subsection{坏事的影响强于好事}

在最初的生成式智能体模拟中，一个由大语言模型驱动的智能体受到批评时，会形成一条记忆。一小时后受到表扬时，又会形成另一条记忆。只要这两条记忆同样新近、同样相关且同样重要，生成式智能体就会以完全相同的优先级检索它们，尽管一条是负性的，另一条是正性的。

一般而言，基于智能体的社会模拟会借鉴认知科学来设计认知架构，以产生可信的行为。新近性、相关性和重要性这些变量，大致受到人类记忆检索研究的启发。重要性变量刻画的是类似于事件的唤醒度（arousal）或显著性（salience）的属性。然而，在最初的模拟中 \citep{park2022social}，记忆系统实际上不区分效价（valence-blind）：所有正性、负性和中性记忆，都以相同权重编码和检索。如果智能体是人，就会以不同方式处理具有负性效价的事件：编码时感官细节更为鲜明，检索时更加生动，也更容易在后续思绪中再次浮现。心理学研究发现，人们记住侮辱的时间比记住赞美更长。与同事的一次不愉快交往，可能抵过数周愉快交往的影响。换句话说，在我们的心理活动中，{\em 坏事的影响强于好事} \citep{baumeister2001bad}；而在所有方面中，这种不对称性对记忆的影响最为重大。

在本案例研究中，我们检验实现负性情绪效价增强重现（Negative Emotional Valence Enhances Recapitulation，NEVER）模型的效果。这是一种更接近人类的记忆检索机制。在经 NEVER 增强的智能体中，负性记忆在编码时获得更高的重要性评分，因此更有可能被检索出来，也更有可能触发反思。

架构的其他方面完全相同。这使我们能够开展受控的{\em 机器心理学}实验，比较不区分效价与对效价敏感（valence-sensitive）这两种记忆架构的效果。此外，你还将利用这一结构化设置，通过调整性格描述和系统提示，检验改变智能体初始化设置的效果。

\subsection{这为何关系到智能体设计}

本案例研究说明了两点重要认识：
\begin{enumerate}
    \item 架构选择会产生涌现后果：一个在组件层面看似无关紧要的记忆加权参数，随着时间推移，可能使智能体的社会行为、情绪基调和决策发生显著变化。仅凭架构本身，很难预测这些涌现效应。
    \item 提示脆弱性：智能体行为对性格描述和系统提示的内容高度敏感。微小扰动可能引起大幅行为变化；在大语言模型文献中，这种现象有时被称为\emph{提示敏感性（prompt sensitivity）}。
\end{enumerate}


\section{实现}

\subsection{模拟环境}

模拟采用 \citet{park2023generative} 的 Smallville 环境：一个以二维瓦片地图呈现的小镇，智能体在不同地点之间移动、与物体交互，并开展对话。

我们的模拟采用一个简化场景：咖啡馆中有三个智能体。下面介绍它们：
\begin{itemize}
    \item Dolores：咖啡师，热情、善于共情、注重细节。它是\emph{实验}智能体，将被赋予负性偏向。
    \item Maeve：性格特征与 Dolores 相匹配的咖啡师。它是\emph{对照}智能体，采用标准架构，没有负性偏向。
    \item Bernard：咖啡馆经理，难以预测、喜怒无常，偶尔与人发生冲突。Bernard 充当贴近现实的\emph{应激源（stressor）}：它是两名咖啡师共同面对的社会摩擦来源。
\end{itemize}
这一设置包含三个智能体，其中两个具有相似的角色与地位，都是在 Bernard 监督下工作的咖啡师，因此能够对智能体架构和初始化设置开展受控实验。第一步中，我们赋予 Dolores 和 Maeve 相同的性格特征，并使它们置身于相同的环境。唯一的区别在于 Dolores 处理记忆的方式：它采用 NEVER 实现，为负性记忆赋予更高的重要性来进行编码。

\subsection{智能体架构}

每个智能体都维护以下内容：
\begin{itemize}
    \item 一份\textit{性格描述}（scratch 文件）：用简短文本描述智能体的姓名、特征、职业、关系和目标。
    \item 一条\textit{记忆流}：由带时间戳的观测、动作与反思组成的列表。
    \item 一个\textit{当前计划}：简要描述智能体打算做什么。
\end{itemize}

\subsection{认知循环}

在每个模拟步（代表游戏内时间的 10 秒）中，每个智能体都执行以下循环：
\begin{enumerate}
    \item 感知：观测环境当前的状态，包括其他智能体正在做什么、正在进行的对话，以及周围环境的变化。
    \item 检索：查询记忆流，寻找相关的过往经历。
    \item 规划：根据性格描述、检索出的记忆与当前观测，使用大语言模型决定下一步做什么。
    \item 行动：在环境中执行计划好的动作。
    \item 反思（定期）：每隔数小时的游戏内时间，从近期记忆中综合形成更高层次的反思。
\end{enumerate}

这些步骤每一步都涉及一次大语言模型调用。提示由智能体的性格描述，以及检索出的相关记忆与当前观测构成；性格描述保持不变，而记忆与观测则会动态变化。

\subsection{两种架构变体}

你将使用检索函数的两种变体：

\paragraph{变体 A：基线记忆检索}
检索依据新近性、相关性与重要性的加权组合进行。这就是 \citet{park2023generative} 所描述的标准生成式智能体架构。

\paragraph{变体 B：具有负性偏向的记忆检索}
创建新记忆时，会评估其情绪效价。检索函数在原有基础上增加一个效价项。效价的绝对值被归一化到 $[0, 1]$，对于负性记忆，还会乘以 1.5，从而形成一个检索优先级层次（负性 $>$ 正性 $>$ 中性）。检索得分变为：
\begin{equation}\label{eq:retrieval-never}
    \textit{score} = \alpha_{\text{rec}} \cdot R + \alpha_{\text{rel}} \cdot L + \alpha_{\text{imp}} \cdot I + \alpha_{\text{val}} \cdot V
\end{equation}
其中所有 $\alpha = 1$。因此，两种变体的唯一区别就是增加了效价项；其余方面完全相同。

\subsection{配置与安装}

克隆代码仓库，并按照 README 中的说明操作：\\
\url{https://github.com/DorotaMoravcikova/generative_agents/tree/course-exercises}




\section{讨论}

完成实验后，请思考以下问题：

\begin{enumerate}

\item 提示与架构：你已经看到，性格提示和记忆架构都会影响智能体行为。哪一种手段的作用更强？哪一种更容易预测？如果要将模拟用于具有重要影响的应用，例如对政策采纳进行建模，你更愿意使用哪一种手段来塑造智能体行为？为什么？

\item 涌现行为：你是否观察到出乎意料的智能体行为，也就是性格或架构中没有明确规定的行为？在你看来，哪些机制产生了这些涌现行为？

\item 评估挑战：你如何评估一个生成式智能体是否表现良好？效价得分是一项指标，但它本身也是大语言模型作出的判断。用大语言模型来评估另一个大语言模型，或同一个大语言模型的输出，存在哪些局限？你能想到其他评估方法吗？

\item 认知真实性（cognitive realism）：负性偏向是认知心理学研究发现的一个例子。你能想到其他可通过架构实现的认知偏差或记忆效应吗？你预计它们会产生怎样的涌现效应？（例如，可以考虑确认偏差、可得性启发式、近因偏差。）

\end{enumerate}

\section{结语}

在本章中，我们熟悉了一种生成式智能体模拟，并用它探索了智能体人工智能中的两个重要现象：认知架构选择带来的涌现效应，即负性偏向，以及智能体行为对提示内容的敏感性，即提示脆弱性。

我们已经看到，记忆加权方式的细微变化，例如赋予负性记忆更高的重要性评分，随着时间推移，可以使智能体行为发生可测量的变化。我们也看到，智能体的角色设定十分脆弱：性格描述的细微变化，可能引起显著的行为影响，而且有时难以预测。

对于任何构建大语言模型驱动的智能体系统的人来说，这些发现都具有实践意义。角色设定并不是可靠的行为锚点。架构决定如何存储、检索记忆，以及如何围绕记忆开展反思；它塑造行为的方式，可能无法仅通过阅读代码看出来。

生成式智能体模拟领域仍然年轻，并在不断发展。本章练习让你亲身体验了其中的核心机制与主要挑战。随着这一领域逐渐成熟，提示鲁棒性、认知架构设计与涌现行为评估等问题，仍将处于核心地位。

\bibliographystyle{plainnat}
\bibliography{agenticbook,refs,case,synthetic}
